\documentclass[11pt,a4paper]{article}
\usepackage[margin=1in]{geometry}
\usepackage[T1]{fontenc}
\usepackage{mathptmx}
\usepackage{microtype}
\usepackage[hidelinks]{hyperref}
\setlength{\parskip}{0.55em}
\setlength{\parindent}{0pt}
\pagestyle{empty}
\newcommand{\wn}{cm$^{-1}$}

\begin{document}

{\large Spectrochimica Acta Part A: Molecular and Biomolecular Spectroscopy}

12 November 2026

Dear Editors,

We submit for consideration as a Full Length Article the manuscript entitled
``A calibrated single-conformer harmonic density functional theory Raman library does
not reach molecular identification of authentic chemical warfare agents.''

\textbf{The question is one of vibrational spectroscopy.} A computed reference library
replaces a measurement with a forward model --- a harmonic frequency calculation, an
empirical frequency scaling, an activity-to-intensity conversion and an assumed lineshape.
How accurate must that chain be before library search succeeds, and what does its residual
error do to the decision to abstain? Computed libraries have been built for monolayer
materials and inorganic solids, but the question has not been answered against a measured
control on authentic material, and it cannot be answered with a simulant, because a Raman
spectrum reports on exactly the molecular skeleton a simulant alters. Chemical warfare
agents are where it is both answerable and consequential: authentic material can be
measured in few laboratories, so the reference library rather than the spectrometer is the
binding constraint.

\textbf{The experiment.} Eleven sealed vials of nerve, vesicant and fourth-generation agents
were measured at 1064\,nm from raw detector counts in an OPCW designated laboratory, and
thirty acquisitions were queried against a measured and a calibrated 54-member computed
library under one pipeline, one native grid and one metric. No calibration parameter is
fitted to identification: the frequency scale comes from band positions, the linewidth from
measured band widths, and the scattering conversion has no free parameter.

\textbf{The result.} The measured library assigns all thirty acquisitions and the computed
library thirteen, or 12 to 19 across all 32 combinations of the forward model's declared
free choices. The measured figure is a within-vial repeatability ceiling, and we say so: on
another laboratory's published band lists the two libraries reach three of five and two of
five, and that ratio is what transfers. What fails is accuracy on the correct molecule, not
separation between candidates --- similarity to the correct entry falls from 0.983 to 0.703
while congener similarities are comparable --- and the deficit concentrates in the
600--850\,\wn{} window built from P--F, P=O, P--O and P--C stretching, at a band-position residual of 0.55 of the linewidth overall.

\textbf{What does not repair it.} Twelve similarity metrics, with the measured library
carried through the identical code as a control, recover at most three of the seventeen
errors. Three escalations of the level of theory change one assignment, on a structure
later shown to be wrong. A conformer search found two computed structures, GF and A-232,
off the global minimum by 4.57 and 9.47\,kcal\,mol$^{-1}$ at the library's own level of
theory; correcting them lowers the closed-set result to thirteen whichever of the two
defensible pre-correction calibrations it is compared with (seventeen or fifteen), and we
report both. Boltzmann averaging over the conformer ensembles of three agents changes no assignment and exposed
a third structural error, which does not change the identification either. Scaling resolved
by mode type, the form the recent literature prefers, returns the same thirteen on a worse
fit.

\textbf{What the computed library does support, and what it does not.} At the class level
it is correct for 26 of 27 acquisitions while the queried molecule is in the library, and
the correct molecule is in the top three of 54 candidates for 22 of 30. With the query's
own entry deleted --- the case a computed reference exists to cover --- class assignment
falls to 19 of 27 and soman is read as a vesicant in all four of its acquisitions. Of six
abstention scores none separates the class-absent condition without error (best area under
the curve 0.90, agent-level 95\% interval 0.76 to 1.00), whereas the same procedure on the
measured library separates it completely. Reliable abstention is therefore bounded by the
accuracy of the forward model, not by the design of the detector.

\textbf{Fit to the journal.} The calculation, the instrument and the cosine are deliberately
standard; what is new is an identification test with no parameter fitted to the reported
performance, used to measure how accurate a forward model must be. The manuscript is the
measured counterpart of Mostafapour et al.\ (\textit{Spectrochim. Acta A} \textbf{344},
126702, 2026), who showed on synthetic spectra that no similarity metric is uniformly
reliable and that past small displacements every metric they tested inverts. Their
principal dataset is a computed vibrational line list perturbed by design; here the
perturbations are supplied by the calculation, the control is a measured library of the
same molecules, and the physical origin is identified --- a residual of 0.55 linewidths,
the regime in which they report metrics to misbehave.

\textbf{Limits stated in the manuscript.} Eleven agents, one instrument, one campaign and
one vial per agent; no non-agent material, so no false-positive rate is claimed; no
orthogonal assay of the materials, and one vial, labeled HN2, was found to contain benzene.
Every headline figure is reported with and without that vial (27 of 27 against 13 of 27
without it). The manuscript contains no synthesis, precursor, quantity or dissemination
information, and a Responsible disclosure statement sets out our assessment of the
deposited data.

The analysis code, preprocessed spectra, calibrated computed library and ORCA inputs and
outputs are deposited at \url{https://github.com/bisu9082/ramanforge} and archived under
DOI \href{https://doi.org/10.5281/zenodo.24107632}{10.5281/zenodo.24107632}.

The work is original, is not under consideration elsewhere, and all authors have read and
approved the submission. Doo-Hee Lee and Seongchan Lee contributed equally. The authors
declare no competing interests.

Sincerely, on behalf of all authors,

\vspace{1.2em}
Ku Kang, Ph.D.\\
Corresponding author\\
CBRN Defence Research Institute, Ministry of National Defense, Republic of Korea\\
E-mail: \href{mailto:bisu9082@gmail.com}{bisu9082@gmail.com}

\end{document}
